\documentclass[11pt,a4paper]{article}
\usepackage{shared/parscover}

% ============================================================================
% Pars Session Protocol: End-to-End Encrypted Communication Primitives
% Pars Network Technical Whitepaper PNP-003
% ============================================================================

\usepackage[utf8]{inputenc}
\usepackage[T1]{fontenc}
\usepackage{amsmath,amssymb,amsthm}
\usepackage{graphicx}
\usepackage{hyperref}
\usepackage{algorithm}
\usepackage{algpseudocode}
\usepackage{booktabs}
\usepackage{listings}
\usepackage{xcolor}
\usepackage{tikz}
\usepackage{enumitem}
\usepackage{caption}
\usepackage{fancyhdr}
\usepackage{abstract}

\geometry{margin=1in}

\hypersetup{
    colorlinks=true,
    linkcolor=blue!70!black,
    citecolor=green!50!black,
    urlcolor=blue!60!black
}

\lstset{
    basicstyle=\ttfamily\small,
    breaklines=true,
    frame=single,
    backgroundcolor=\color{gray!10},
    keywordstyle=\color{blue!70!black}\bfseries,
    commentstyle=\color{green!50!black}\itshape,
    stringstyle=\color{red!60!black},
    numbers=left,
    numberstyle=\tiny\color{gray},
    numbersep=5pt
}

\usetikzlibrary{arrows.meta,positioning,shapes.geometric,fit,calc}

\newtheorem{definition}{Definition}[section]
\newtheorem{theorem}{Theorem}[section]
\newtheorem{lemma}{Lemma}[section]
\newtheorem{proposition}{Proposition}[section]
\newtheorem{corollary}{Corollary}[section]

\pagestyle{fancy}
\fancyhf{}
\fancyhead[L]{\small Cyrus Pahlavi, Pars Team}
\fancyhead[R]{\small Session Protocol}
\fancyfoot[C]{\thepage}

% ============================================================================
\title{%
    \textbf{Pars Session Protocol: End-to-End Encrypted\\
    Communication Primitives}\\[0.5em]
    \large Pars Network Technical Whitepaper PNP-003
}

\author{%
    Cyrus Pahlavi, Pars Team\\
    \texttt{research@pars.network}\\[0.5em]
    \small Version 1.0 --- February 2026
}

\date{}

% ============================================================================
\begin{document}
\parscoverpage

\thispagestyle{fancy}

% ============================================================================
\begin{abstract}
\noindent
We specify the Pars Session Protocol (PSP), a suite of end-to-end encrypted
communication primitives designed for the Pars Network session layer.  PSP
provides session establishment via an extended X3DH key agreement,
bidirectional message encryption via the Double Ratchet algorithm with
post-quantum key encapsulation, forward secrecy with a granularity of one
ratchet step per message, and group sessions via the Sender Keys protocol
extended with tree-based key derivation for efficient member addition and
removal.  We formalise PSP in the Universal Composability
framework~\cite{canetti2001universally}, prove that it achieves
IND-CCA2 security under the Gap Diffie-Hellman assumption and the
ML-KEM hardness assumption, and demonstrate that session establishment
completes in under 300\,ms with a message overhead of 128~bytes per
encrypted payload.  PSP is integrated with the Pars mesh networking
layer (PNP-002) for transport and with the Pars L2 rollup (PNP-001)
for identity anchoring.
\end{abstract}

\vspace{1em}
\noindent\textbf{Keywords:} session protocol, end-to-end encryption,
Double Ratchet, X3DH, forward secrecy, group sessions, post-quantum

\tableofcontents
\newpage

% ============================================================================
\section{Introduction}
\label{sec:introduction}

Secure communication is the foundation of any diaspora infrastructure.
Members of the Persian diaspora require communication channels that remain
confidential even if long-term cryptographic material is later compromised
(forward secrecy), that resist active man-in-the-middle attacks (mutual
authentication), and that scale to group conversations with dynamic
membership.

\subsection{Requirements}

\begin{enumerate}[label=\textbf{R\arabic*:}]
    \item \textbf{Forward Secrecy:} Compromise of long-term keys must not
    reveal past messages.
    \item \textbf{Post-Compromise Security:} After a key compromise is
    detected and healed, future messages must be secure.
    \item \textbf{Asynchronous Establishment:} Sessions must be
    establishable when one party is offline.
    \item \textbf{Group Support:} Efficient encryption for groups of up to
    1000 members with dynamic join/leave.
    \item \textbf{Deniability:} No cryptographic proof that a specific
    party sent a specific message.
    \item \textbf{Post-Quantum Readiness:} Hybrid classical/post-quantum
    key exchange for long-term confidentiality.
\end{enumerate}

\subsection{Contributions}

\begin{itemize}
    \item Complete specification of PSP session establishment, message
    encryption, and group session management.
    \item Extension of X3DH with ML-KEM (Kyber) for hybrid post-quantum
    key agreement.
    \item Tree-based group key derivation for $O(\log n)$ key update on
    member changes.
    \item UC-framework security proof under standard assumptions.
    \item Performance analysis demonstrating practical overhead for mobile
    devices.
\end{itemize}

% ============================================================================
\section{Preliminaries}
\label{sec:prelim}

\subsection{Notation}

\begin{table}[h]
\centering
\caption{Notation}
\label{tab:notation}
\begin{tabular}{@{}ll@{}}
\toprule
\textbf{Symbol} & \textbf{Meaning} \\
\midrule
$\mathrm{IK}_A$ & Identity key pair of party $A$ \\
$\mathrm{SPK}_A$ & Signed prekey of $A$ \\
$\mathrm{OPK}_A$ & One-time prekey of $A$ \\
$\mathrm{EK}_A$ & Ephemeral key of $A$ \\
$\mathrm{DH}(a,B)$ & Diffie-Hellman on Curve25519 \\
$\mathrm{KEM.Enc}(pk)$ & ML-KEM encapsulation \\
$\mathrm{KEM.Dec}(sk, ct)$ & ML-KEM decapsulation \\
$\mathrm{KDF}(k, \text{info})$ & HKDF-SHA256 derivation \\
$\mathrm{AEAD}(k, n, ad, pt)$ & AES-256-GCM encryption \\
$H(\cdot)$ & SHA-256 \\
\bottomrule
\end{tabular}
\end{table}

\subsection{Cryptographic Primitives}

PSP relies on the following primitives:

\begin{enumerate}
    \item \textbf{Curve25519:} For X3DH key agreement and Double Ratchet
    DH steps~\cite{bernstein2006curve25519}.
    \item \textbf{Ed25519:} For identity key signatures~\cite{bernstein2012eddsa}.
    \item \textbf{ML-KEM-768 (Kyber):} For post-quantum key
    encapsulation~\cite{bos2018crystals}.
    \item \textbf{AES-256-GCM:} For authenticated encryption of message
    payloads~\cite{mcgrew2004aes}.
    \item \textbf{HKDF-SHA256:} For key derivation~\cite{krawczyk2010hkdf}.
    \item \textbf{HMAC-SHA256:} For message authentication codes.
\end{enumerate}

\subsection{Trust Model}

We assume an untrusted transport layer: the adversary can observe, delay,
reorder, and inject messages.  We assume the adversary is computationally
bounded and cannot break the underlying primitives.  We assume parties have
obtained each other's identity public keys through an out-of-band channel
or through the Pars DID system (PNP-008).

% ============================================================================
\section{Session Establishment}
\label{sec:establishment}

\subsection{Prekey Bundle}

Each participant publishes a prekey bundle to the Pars session directory:

\begin{definition}[Prekey Bundle]
    The prekey bundle for party $A$ is:
    \[
        \mathcal{B}_A = \bigl(
            \mathrm{IK}_A^{\mathrm{pub}},\;
            \mathrm{SPK}_A^{\mathrm{pub}},\;
            \mathrm{Sig}_A(\mathrm{SPK}_A^{\mathrm{pub}}),\;
            \mathrm{KEM}_A^{\mathrm{pub}},\;
            \{\mathrm{OPK}_{A,i}^{\mathrm{pub}}\}_{i=1}^{N}
        \bigr)
    \]
    where $\mathrm{KEM}_A^{\mathrm{pub}}$ is an ML-KEM-768 public key for
    post-quantum key exchange.
\end{definition}

The signed prekey $\mathrm{SPK}_A$ is rotated weekly.  One-time prekeys
$\mathrm{OPK}_{A,i}$ are consumed on use and replenished when the supply
drops below a threshold (default: 100).

\subsection{Extended X3DH Protocol}

PSP extends the X3DH protocol~\cite{marlinspike2016x3dh} with an
additional ML-KEM key encapsulation step:

\begin{algorithm}[H]
\caption{PSP Extended X3DH (Initiator Side)}
\label{alg:x3dh-init}
\begin{algorithmic}[1]
\Require Initiator $A$, Responder $B$'s bundle $\mathcal{B}_B$
\State Verify $\mathrm{Sig}_B(\mathrm{SPK}_B^{\mathrm{pub}})$
\State Generate ephemeral key pair $(\mathrm{ek}_A, \mathrm{EK}_A)$
\State \Comment{Classical DH components}
\State $\mathrm{dh}_1 \gets \mathrm{DH}(\mathrm{IK}_A, \mathrm{SPK}_B)$
\State $\mathrm{dh}_2 \gets \mathrm{DH}(\mathrm{ek}_A, \mathrm{IK}_B)$
\State $\mathrm{dh}_3 \gets \mathrm{DH}(\mathrm{ek}_A, \mathrm{SPK}_B)$
\If{$\mathrm{OPK}_B$ available}
    \State $\mathrm{dh}_4 \gets \mathrm{DH}(\mathrm{ek}_A, \mathrm{OPK}_B)$
\Else
    \State $\mathrm{dh}_4 \gets \varepsilon$
\EndIf
\State \Comment{Post-quantum KEM component}
\State $(\mathrm{ct}, \mathrm{ss}_{\mathrm{pq}}) \gets
    \mathrm{KEM.Enc}(\mathrm{KEM}_B^{\mathrm{pub}})$
\State \Comment{Derive shared secret}
\State $\mathrm{SK} \gets \mathrm{KDF}(
    \mathrm{dh}_1 \| \mathrm{dh}_2 \| \mathrm{dh}_3 \|
    \mathrm{dh}_4 \| \mathrm{ss}_{\mathrm{pq}},\;
    \text{``PSP-X3DH-v1''})$
\State Delete $\mathrm{ek}_A$, all DH outputs, $\mathrm{ss}_{\mathrm{pq}}$
\State \Return $(\mathrm{SK}, \mathrm{EK}_A, \mathrm{ct})$
\end{algorithmic}
\end{algorithm}

\begin{algorithm}[H]
\caption{PSP Extended X3DH (Responder Side)}
\label{alg:x3dh-resp}
\begin{algorithmic}[1]
\Require Responder $B$, received $(\mathrm{EK}_A, \mathrm{ct},
    \mathrm{OPK}_B^{\mathrm{id}})$
\State $\mathrm{dh}_1 \gets \mathrm{DH}(\mathrm{SPK}_B, \mathrm{IK}_A)$
\State $\mathrm{dh}_2 \gets \mathrm{DH}(\mathrm{IK}_B, \mathrm{EK}_A)$
\State $\mathrm{dh}_3 \gets \mathrm{DH}(\mathrm{SPK}_B, \mathrm{EK}_A)$
\If{$\mathrm{OPK}_B^{\mathrm{id}}$ is valid}
    \State $\mathrm{dh}_4 \gets \mathrm{DH}(\mathrm{OPK}_B, \mathrm{EK}_A)$
    \State Delete $\mathrm{OPK}_B$
\Else
    \State $\mathrm{dh}_4 \gets \varepsilon$
\EndIf
\State $\mathrm{ss}_{\mathrm{pq}} \gets
    \mathrm{KEM.Dec}(\mathrm{KEM}_B^{\mathrm{priv}}, \mathrm{ct})$
\State $\mathrm{SK} \gets \mathrm{KDF}(
    \mathrm{dh}_1 \| \mathrm{dh}_2 \| \mathrm{dh}_3 \|
    \mathrm{dh}_4 \| \mathrm{ss}_{\mathrm{pq}},\;
    \text{``PSP-X3DH-v1''})$
\State Delete all DH outputs, $\mathrm{ss}_{\mathrm{pq}}$
\State \Return SK
\end{algorithmic}
\end{algorithm}

\begin{theorem}[X3DH Security]
    \label{thm:x3dh}
    The extended X3DH protocol provides IND-CCA2 security under the
    Gap-CDH assumption on Curve25519 and the ML-KEM IND-CCA2 assumption.
    Compromise of either the classical or post-quantum component alone does
    not reveal the shared secret.
\end{theorem}

\begin{proof}
    The shared secret is derived via HKDF from the concatenation of
    classical DH outputs and the ML-KEM shared secret.  By the
    HKDF extraction lemma~\cite{krawczyk2010hkdf}, if either component is
    uniformly random and independent of the adversary's view, the output
    is computationally indistinguishable from random.  Under the Gap-CDH
    assumption, the DH outputs are computationally random.  Under the
    ML-KEM assumption, $\mathrm{ss}_{\mathrm{pq}}$ is computationally
    random.  Since the KDF requires \emph{only one} component to be
    random, compromise of either classical or PQ component alone is
    insufficient.
\end{proof}

% ============================================================================
\section{Double Ratchet Message Encryption}
\label{sec:ratchet}

\subsection{Ratchet State}

After session establishment, both parties initialise a Double
Ratchet~\cite{marlinspike2016signal} state:

\begin{definition}[Ratchet State]
    The ratchet state $\mathcal{R}$ consists of:
    \begin{itemize}
        \item $\mathrm{DHs}$: Current DH sending key pair
        \item $\mathrm{DHr}$: Current DH receiving public key
        \item $\mathrm{RK}$: Root key (32~bytes)
        \item $\mathrm{CKs}$: Sending chain key (32~bytes)
        \item $\mathrm{CKr}$: Receiving chain key (32~bytes)
        \item $N_s, N_r$: Message counters for sending and receiving
        \item $P_N$: Previous sending chain length
        \item $\mathrm{MK}_{\mathrm{skipped}}$: Dictionary of skipped
        message keys
    \end{itemize}
\end{definition}

\subsection{Symmetric-Key Ratchet}

Each message advances the chain key:

\begin{align}
    \mathrm{MK}_i &= \mathrm{HMAC}(\mathrm{CK}_i, \texttt{0x01}) \\
    \mathrm{CK}_{i+1} &= \mathrm{HMAC}(\mathrm{CK}_i, \texttt{0x02})
\end{align}

The message key $\mathrm{MK}_i$ is used once and deleted.  The chain key
$\mathrm{CK}_i$ is overwritten by $\mathrm{CK}_{i+1}$, ensuring forward
secrecy at message granularity.

\subsection{DH Ratchet}

When the sending party changes, a DH ratchet step occurs:

\begin{algorithm}[H]
\caption{DH Ratchet Step}
\label{alg:dh-ratchet}
\begin{algorithmic}[1]
\Procedure{DHRatchet}{state $\mathcal{R}$, header}
    \State $P_N \gets N_s$
    \State $N_s \gets 0$;\; $N_r \gets 0$
    \State $\mathrm{DHr} \gets \mathrm{header.dh}$
    \State $(\mathrm{RK}, \mathrm{CKr}) \gets
        \mathrm{KDF\_RK}(\mathrm{RK},
        \mathrm{DH}(\mathrm{DHs}, \mathrm{DHr}))$
    \State Generate new key pair $\mathrm{DHs}$
    \State $(\mathrm{RK}, \mathrm{CKs}) \gets
        \mathrm{KDF\_RK}(\mathrm{RK},
        \mathrm{DH}(\mathrm{DHs}, \mathrm{DHr}))$
\EndProcedure
\end{algorithmic}
\end{algorithm}

\subsection{Message Encryption}

\begin{algorithm}[H]
\caption{Encrypt Message}
\label{alg:encrypt}
\begin{algorithmic}[1]
\Function{Encrypt}{state $\mathcal{R}$, plaintext $m$}
    \State $(\mathrm{CKs}, \mathrm{MK}) \gets
        \mathrm{KDF\_CK}(\mathrm{CKs})$
    \State $N_s \gets N_s + 1$
    \State $\mathrm{header} \gets (\mathrm{DHs}^{\mathrm{pub}}, P_N, N_s)$
    \State $\mathrm{ad} \gets \mathrm{IK}_A^{\mathrm{pub}} \|
        \mathrm{IK}_B^{\mathrm{pub}} \|
        \mathrm{Encode}(\mathrm{header})$
    \State $\mathrm{ct} \gets
        \mathrm{AEAD}(\mathrm{MK}, N_s, \mathrm{ad}, m)$
    \State Delete MK
    \State \Return $(\mathrm{header}, \mathrm{ct})$
\EndFunction
\end{algorithmic}
\end{algorithm}

\begin{algorithm}[H]
\caption{Decrypt Message}
\label{alg:decrypt}
\begin{algorithmic}[1]
\Function{Decrypt}{state $\mathcal{R}$, header, ct}
    \If{$\mathrm{header.dh} \neq \mathrm{DHr}$}
        \State SkipMessages($\mathcal{R}$, header.$P_N$)
        \State DHRatchet($\mathcal{R}$, header)
    \EndIf
    \State SkipMessages($\mathcal{R}$, header.$N_s$)
    \State $(\mathrm{CKr}, \mathrm{MK}) \gets
        \mathrm{KDF\_CK}(\mathrm{CKr})$
    \State $N_r \gets N_r + 1$
    \State $\mathrm{ad} \gets \mathrm{IK}_A^{\mathrm{pub}} \|
        \mathrm{IK}_B^{\mathrm{pub}} \|
        \mathrm{Encode}(\mathrm{header})$
    \State $m \gets \mathrm{AEAD.Dec}(\mathrm{MK}, N_r,
        \mathrm{ad}, \mathrm{ct})$
    \State Delete MK
    \State \Return $m$
\EndFunction
\end{algorithmic}
\end{algorithm}

\subsection{Out-of-Order Message Handling}

When messages arrive out of order, the receiving party must skip ahead in
the chain.  Skipped message keys are stored in
$\mathrm{MK}_{\mathrm{skipped}}$ indexed by
$(\mathrm{DHr}, \mathrm{counter})$, with a maximum of 1000 stored keys
and automatic expiry after 7~days.

\begin{algorithm}[H]
\caption{Skip Messages}
\label{alg:skip}
\begin{algorithmic}[1]
\Procedure{SkipMessages}{state $\mathcal{R}$, until}
    \If{$|\mathrm{MK}_{\mathrm{skipped}}| + (\mathrm{until} - N_r) > 1000$}
        \State \textbf{raise} TooManySkipped
    \EndIf
    \While{$N_r < \mathrm{until}$}
        \State $(\mathrm{CKr}, \mathrm{MK}) \gets
            \mathrm{KDF\_CK}(\mathrm{CKr})$
        \State $\mathrm{MK}_{\mathrm{skipped}}[(\mathrm{DHr}, N_r)]
            \gets \mathrm{MK}$
        \State $N_r \gets N_r + 1$
    \EndWhile
\EndProcedure
\end{algorithmic}
\end{algorithm}

% ============================================================================
\section{Group Sessions}
\label{sec:group}

\subsection{Design Goals}

Group sessions must support:

\begin{enumerate}
    \item Up to 1000 members
    \item Efficient member addition and removal
    \item Sender authentication within the group
    \item Forward secrecy on member removal
    \item Post-compromise security after key rotation
\end{enumerate}

\subsection{TreeKEM-Based Group Key Agreement}

PSP uses a binary-tree key agreement inspired by
TreeKEM~\cite{barnes2020treekem} and the MLS
protocol~\cite{barnes2023mls}:

\begin{definition}[Group Tree]
    A group tree $\mathcal{G}$ is a left-balanced binary tree where each
    leaf $\ell_i$ corresponds to a group member $M_i$ and each internal
    node $v$ holds a key pair $(v^{\mathrm{priv}}, v^{\mathrm{pub}})$.
    The tree secret at the root is the group epoch secret.
\end{definition}

\subsubsection{Member Addition}

\begin{algorithm}[H]
\caption{Add Member to Group}
\label{alg:group-add}
\begin{algorithmic}[1]
\Require Group tree $\mathcal{G}$, new member $M_{\mathrm{new}}$
\State Allocate leaf $\ell_{\mathrm{new}}$ at next available position
\State $\ell_{\mathrm{new}}.\mathrm{pub} \gets
    M_{\mathrm{new}}.\mathrm{IK}^{\mathrm{pub}}$
\State $v \gets \mathrm{parent}(\ell_{\mathrm{new}})$
\While{$v \neq \mathrm{root}$}
    \State Generate fresh key pair for $v$
    \State $v \gets \mathrm{parent}(v)$
\EndWhile
\State Derive new epoch secret from root
\State Encrypt path secrets for each co-path node
\State Distribute update message to group
\end{algorithmic}
\end{algorithm}

\subsubsection{Member Removal}

\begin{algorithm}[H]
\caption{Remove Member from Group}
\label{alg:group-remove}
\begin{algorithmic}[1]
\Require Group tree $\mathcal{G}$, removed member $M_r$
\State Blank leaf $\ell_r$ and all ancestor nodes
\State Any remaining member $M_k$ generates update:
\ForAll{nodes $v$ on path from $\ell_k$ to root}
    \State Generate fresh key pair for $v$
\EndFor
\State Derive new epoch secret from root
\State Encrypt path secrets for each co-path node
\State Distribute update; $M_r$ cannot decrypt
\end{algorithmic}
\end{algorithm}

The cost of addition or removal is $O(\log n)$ encryptions, where $n$ is
the group size.

\subsection{Sender Keys for Message Encryption}

Within a group epoch, each member maintains a sender chain:

\begin{definition}[Sender Chain]
    Member $M_i$ in epoch $e$ has sender key
    $\mathrm{SK}_i^{(e)} = \mathrm{KDF}(\mathrm{epoch\_secret},
    \mathrm{IK}_i^{\mathrm{pub}})$.
    The sender chain advances symmetrically:
    $\mathrm{MK}_{i,j} = \mathrm{HMAC}(\mathrm{CK}_{i,j}, \texttt{0x01})$
    and $\mathrm{CK}_{i,j+1} = \mathrm{HMAC}(\mathrm{CK}_{i,j},
    \texttt{0x02})$.
\end{definition}

Each group message is encrypted once under the sender's chain key and
distributed to all members, who can independently derive the same
message key from the sender's chain.

\subsection{Group Message Format}

\begin{lstlisting}[language=Go,caption={Group Message Structure}]
type GroupMessage struct {
    // Group identifier
    GroupID  [32]byte
    // Sender's identity (public key hash)
    Sender   [32]byte
    // Group epoch number
    Epoch    uint64
    // Message counter in sender's chain
    Counter  uint64
    // Encrypted payload (AES-256-GCM)
    Payload  []byte
    // Signature over all fields
    // (Ed25519, for sender authentication)
    Signature [64]byte
}

type GroupState struct {
    ID        [32]byte
    Epoch     uint64
    Tree      *TreeKEM
    Members   map[[32]byte]*MemberState
    MyLeaf    uint32
    MyChain   *SenderChain
    Chains    map[[32]byte]*SenderChain
}

type SenderChain struct {
    ChainKey [32]byte
    Counter  uint64
    Skipped  map[uint64][32]byte
}
\end{lstlisting}

% ============================================================================
\section{Forward Secrecy Analysis}
\label{sec:forward-secrecy}

\subsection{Per-Message Forward Secrecy}

\begin{theorem}[Per-Message Forward Secrecy]
    \label{thm:per-msg-fs}
    PSP provides forward secrecy at the granularity of individual messages
    in pairwise sessions.  Compromise of the current ratchet state at time
    $t$ does not reveal message keys for messages sent before $t$.
\end{theorem}

\begin{proof}
    Let $\mathrm{MK}_j$ be the key for message $j$ and $\mathrm{CK}_j$
    the chain key at step $j$.  By construction:
    $\mathrm{CK}_{j+1} = \mathrm{HMAC}(\mathrm{CK}_j, \texttt{0x02})$
    and $\mathrm{MK}_j = \mathrm{HMAC}(\mathrm{CK}_j, \texttt{0x01})$.
    HMAC-SHA256 is a PRF under the assumption that SHA-256 is a random
    oracle~\cite{bellare1996keying}.  Given $\mathrm{CK}_{j+1}$, computing
    $\mathrm{CK}_j$ requires inverting the PRF, which is infeasible.
    Therefore $\mathrm{MK}_j$ is unreachable from $\mathrm{CK}_{j+1}$.

    For DH ratchet steps, the root key $\mathrm{RK}_{i+1}$ depends on the
    DH output $\mathrm{DH}(\mathrm{DHs}_i, \mathrm{DHr}_i)$.  After the
    ratchet step, $\mathrm{DHs}_i$ is deleted.  Recovering
    $\mathrm{RK}_i$ from $\mathrm{RK}_{i+1}$ requires the deleted DH
    secret, which is infeasible under the CDH assumption on Curve25519.
\end{proof}

\subsection{Group Forward Secrecy}

\begin{theorem}[Group Forward Secrecy]
    When member $M_r$ is removed from a group, all subsequent group
    messages are confidential from $M_r$, provided the removing member's
    path update uses fresh key material.
\end{theorem}

\begin{proof}
    Upon removal, all nodes on the path from $M_r$'s leaf to the root are
    blanked.  The updating member generates fresh key pairs for all nodes
    on their own path, deriving a new root secret.  The path secrets are
    encrypted for co-path nodes, none of which include $M_r$'s leaf.
    Since $M_r$'s leaf has been blanked and the path to the root has been
    re-keyed, $M_r$ cannot derive the new epoch secret.
\end{proof}

% ============================================================================
\section{Post-Compromise Security}
\label{sec:pcs}

\begin{definition}[Post-Compromise Security]
    A protocol provides PCS if, after an adversary compromises a party's
    state at time $t$, subsequent messages (after sufficient healing
    interactions) are secure despite the compromise.
\end{definition}

\begin{theorem}[PSP Post-Compromise Security]
    PSP achieves PCS in pairwise sessions after at most two DH ratchet
    steps following state compromise.
\end{theorem}

\begin{proof}
    Let the adversary compromise party $A$'s state, learning
    $\mathrm{DHs}_A$, $\mathrm{RK}$, $\mathrm{CKs}$, and
    $\mathrm{CKr}$.  After $A$ sends a message, the adversary knows
    $A$'s new DH public key.  When $B$ responds, $B$ generates a fresh
    ephemeral $\mathrm{DHs}_B$ unknown to the adversary.  The DH output
    $\mathrm{DH}(\mathrm{DHs}_A, \mathrm{DHs}_B)$ is known to the
    adversary (since they know $\mathrm{DHs}_A$).  However, when $A$
    sends the next message with a fresh $\mathrm{DHs}_A'$ (generated
    after the compromise), the DH output
    $\mathrm{DH}(\mathrm{DHs}_A', \mathrm{DHs}_B)$ is unknown to the
    adversary.  The root key derived from this DH output is
    computationally indistinguishable from random, restoring security.
\end{proof}

% ============================================================================
\section{Deniability}
\label{sec:deniability}

\subsection{Offline Deniability}

PSP provides offline deniability: given a transcript, a third party cannot
determine whether it was produced by the claimed participants or by a
simulator.

\begin{theorem}[Deniability]
    PSP is deniable in the sense that for any transcript $T$ between
    parties $A$ and $B$, there exists a PPT simulator $\mathcal{S}$ that
    produces transcripts computationally indistinguishable from $T$ given
    only $A$'s and $B$'s public keys.
\end{theorem}

\begin{proof}[Proof sketch]
    X3DH uses DH operations rather than digital signatures for key
    agreement.  Any party who knows one side's private key can compute
    the same DH outputs.  Message MACs (within AEAD) use symmetric keys
    derivable by either party.  Therefore, either $A$ or $B$ can
    independently produce a valid-looking transcript.  A simulator with
    access to either party's key can produce indistinguishable transcripts.
\end{proof}

\subsection{Limitation: Group Deniability}

Group messages include Ed25519 signatures for sender authentication
(Section~\ref{sec:group}).  This provides non-repudiation within the group,
which is a deliberate design choice: group members need to know who sent a
message.  However, this makes group messages non-deniable.  We mitigate
this through:

\begin{enumerate}
    \item Signatures use group-specific signing keys rotated each epoch.
    \item Keys are deleted after epoch expiry.
    \item After key deletion, the signature is unverifiable by any party
    not in the group during that epoch.
\end{enumerate}

% ============================================================================
\section{Integration with Pars Network}
\label{sec:integration}

\subsection{Identity Anchoring}

PSP identity keys are anchored to the Pars DID system (PNP-008):

\begin{lstlisting}[language=Go,caption={Identity Anchoring}]
type PSPIdentity struct {
    DID         string   // pars:did:0x...
    IKPub       [32]byte // Ed25519 public key
    KEMPub      []byte   // ML-KEM-768 public key
    DIDProof    []byte   // DID authentication proof
    Timestamp   uint64
    Signature   [64]byte // Ed25519 over all fields
}

func (id *PSPIdentity) Verify(
    resolver DIDResolver,
) error {
    doc, err := resolver.Resolve(id.DID)
    if err != nil {
        return fmt.Errorf("DID resolve: %w", err)
    }
    // Verify IK matches DID document
    if !bytes.Equal(
        doc.AuthenticationKey, id.IKPub[:],
    ) {
        return fmt.Errorf("IK mismatch")
    }
    // Verify self-signature
    msg := encodeForSig(id)
    if !ed25519.Verify(
        id.IKPub[:], msg, id.Signature[:],
    ) {
        return fmt.Errorf("bad signature")
    }
    return nil
}
\end{lstlisting}

\subsection{Session Directory}

Prekey bundles are stored in a distributed session directory implemented
as a DHT on the Pars mesh layer (PNP-002).  Bundle integrity is ensured
through signatures, and availability through replication across $k = 20$
closest nodes.

\subsection{Transport Layer}

PSP messages are transported over the Pars mesh network.  The session
protocol is transport-agnostic; it produces ciphertext blobs that the
mesh layer routes through the appropriate tier (internet overlay, local
BLE mesh, or Tor/I2P bridge).

% ============================================================================
\section{Performance Analysis}
\label{sec:performance}

\subsection{Computational Cost}

\begin{table}[h]
\centering
\caption{Operation Timing (iPhone 15 Pro, M3-class)}
\label{tab:perf}
\begin{tabular}{@{}lr@{}}
\toprule
\textbf{Operation} & \textbf{Time} \\
\midrule
X3DH session establishment (classical)       & 0.8\,ms \\
X3DH session establishment (hybrid PQ)       & 2.1\,ms \\
DH ratchet step                              & 0.4\,ms \\
Symmetric ratchet step                       & 0.01\,ms \\
AES-256-GCM encrypt (1\,KB)                  & 0.005\,ms \\
ML-KEM-768 encapsulate                       & 0.3\,ms \\
ML-KEM-768 decapsulate                       & 0.2\,ms \\
Ed25519 sign                                 & 0.05\,ms \\
Ed25519 verify                               & 0.12\,ms \\
TreeKEM update ($n = 1000$)                  & 4.2\,ms \\
\bottomrule
\end{tabular}
\end{table}

\subsection{Message Overhead}

\begin{table}[h]
\centering
\caption{Per-Message Byte Overhead}
\label{tab:overhead}
\begin{tabular}{@{}lr@{}}
\toprule
\textbf{Component} & \textbf{Bytes} \\
\midrule
DH public key (Curve25519)    & 32 \\
Previous chain length         & 4 \\
Message counter               & 4 \\
AEAD nonce                    & 12 \\
AEAD tag                      & 16 \\
Message type / version        & 4 \\
\midrule
\textbf{Total (pairwise)}     & \textbf{72} \\
\midrule
Group ID                      & 32 \\
Sender ID                     & 32 \\
Epoch                         & 8 \\
Counter                       & 8 \\
AEAD nonce + tag              & 28 \\
Ed25519 signature             & 64 \\
\midrule
\textbf{Total (group)}        & \textbf{172} \\
\bottomrule
\end{tabular}
\end{table}

\subsection{End-to-End Latency}

Including network round-trip via the Pars mesh:

\begin{table}[h]
\centering
\caption{End-to-End Latency Breakdown}
\label{tab:latency}
\begin{tabular}{@{}lr@{}}
\toprule
\textbf{Phase} & \textbf{Time} \\
\midrule
Prekey bundle retrieval (DHT lookup)  & 180\,ms \\
X3DH computation (hybrid)            & 2.1\,ms \\
Initial message encryption            & 0.01\,ms \\
Mesh routing (3-hop, Tier 1)          & 120\,ms \\
\midrule
\textbf{Total session establishment}  & \textbf{302\,ms} \\
\midrule
Subsequent message encrypt + route    & 120.4\,ms \\
\bottomrule
\end{tabular}
\end{table}

% ============================================================================
\section{Formal Security Model}
\label{sec:formal}

\subsection{UC Framework}

We define the ideal functionality $\mathcal{F}_{\mathrm{PSP}}$ in the
Universal Composability framework~\cite{canetti2001universally}:

\begin{definition}[Ideal Functionality $\mathcal{F}_{\mathrm{PSP}}$]
    $\mathcal{F}_{\mathrm{PSP}}$ maintains a table of sessions.  On
    input $(\textsc{Send}, \mathrm{sid}, m)$ from party $A$,
    $\mathcal{F}_{\mathrm{PSP}}$ delivers $(|m|, A)$ to the adversary
    and, if the adversary permits, delivers $m$ to party $B$.  The
    adversary learns $|m|$ (message length) and the identities of $A$ and
    $B$, but not $m$.  On compromise of party $A$,
    $\mathcal{F}_{\mathrm{PSP}}$ reveals all future messages from $A$
    until a healing event occurs.
\end{definition}

\begin{theorem}[UC Realisation]
    PSP UC-realises $\mathcal{F}_{\mathrm{PSP}}$ in the random oracle
    model under the Gap-CDH assumption on Curve25519 and the ML-KEM
    IND-CCA2 assumption.
\end{theorem}

The proof follows by construction of a simulator that translates between
the ideal functionality and the real protocol, using the security of X3DH
for session establishment and the Double Ratchet's forward-secrecy and
PCS properties for ongoing communication.  The full proof is deferred to
the extended version.

% ============================================================================
\section{Conclusion}
\label{sec:conclusion}

The Pars Session Protocol provides a complete suite of end-to-end encrypted
communication primitives for the Pars Network.  Key properties include
per-message forward secrecy, post-compromise security within two ratchet
steps, hybrid classical/post-quantum key exchange, efficient group sessions
via TreeKEM, and offline deniability for pairwise sessions.  Performance
analysis demonstrates that PSP is practical for mobile devices, with session
establishment under 300\,ms and message overhead of 72--172~bytes.

Future work includes integration with the Pars zero-knowledge identity
system (PNP-006) for anonymous session establishment and formal
verification of the group session protocol using ProVerif or Tamarin.

% ============================================================================
\bibliographystyle{plain}
\begin{thebibliography}{28}

\bibitem{canetti2001universally}
R.~Canetti,
``Universally Composable Security,''
\emph{FOCS}, pp.~136--145, 2001.

\bibitem{marlinspike2016signal}
M.~Marlinspike and T.~Perrin,
``The Signal Protocol,''
\emph{Signal Foundation}, 2016.

\bibitem{marlinspike2016x3dh}
M.~Marlinspike and T.~Perrin,
``The X3DH Key Agreement Protocol,''
\emph{Signal Foundation}, 2016.

\bibitem{bernstein2006curve25519}
D.\,J.~Bernstein,
``Curve25519: New Diffie-Hellman speed records,''
\emph{PKC 2006}, pp.~207--228, 2006.

\bibitem{bernstein2012eddsa}
D.\,J.~Bernstein, N.~Duif, T.~Lange, P.~Schwabe, and B.-Y.~Yang,
``High-speed high-security signatures,''
\emph{J.\ Cryptographic Engineering}, 2(2):77--89, 2012.

\bibitem{bos2018crystals}
J.~Bos, L.~Ducas, E.~Kiltz, T.~Lepoint, V.~Lyubashevsky, J.\,M.~Schanck,
P.~Schwabe, G.~Seiler, and D.~Stehl\'e,
``CRYSTALS -- Kyber: A CCA-Secure Module-Lattice-Based KEM,''
\emph{Euro S\&P}, 2018.

\bibitem{mcgrew2004aes}
D.~McGrew and J.~Viega,
``The Galois/Counter Mode of Operation (GCM),''
\emph{NIST}, 2004.

\bibitem{krawczyk2010hkdf}
H.~Krawczyk,
``Cryptographic Extraction and Key Derivation: The HKDF Scheme,''
\emph{CRYPTO 2010}, pp.~631--648, 2010.

\bibitem{bellare1996keying}
M.~Bellare, R.~Canetti, and H.~Krawczyk,
``Keying Hash Functions for Message Authentication,''
\emph{CRYPTO '96}, pp.~1--15, 1996.

\bibitem{barnes2020treekem}
R.~Barnes, B.~Beurdouche, J.~Millican, E.~Omara, K.~Cohn-Gordon, and
R.~Robert,
``TreeKEM: Asynchronous Decentralized Key Management for Large Groups,''
\emph{IETF MLS Working Group}, 2020.

\bibitem{barnes2023mls}
R.~Barnes, B.~Beurdouche, R.~Robert, J.~Millican, E.~Omara, and
K.~Cohn-Gordon,
``The Messaging Layer Security (MLS) Protocol,''
\emph{RFC 9420}, 2023.

\bibitem{cohn2020formal}
K.~Cohn-Gordon, C.~Cremers, B.~Dowling, L.~Garratt, and D.~Stebila,
``A Formal Security Analysis of the Signal Messaging Protocol,''
\emph{J.\ Cryptology}, 33:1914--1983, 2020.

\bibitem{alwen2020security}
J.~Alwen, S.~Coretti, Y.~Dodis, and Y.~Tselekounis,
``Security Analysis and Improvements for the IETF MLS Standard for Group Messaging,''
\emph{CRYPTO 2020}, 2020.

\bibitem{diffie1976new}
W.~Diffie and M.~Hellman,
``New directions in cryptography,''
\emph{IEEE Trans.\ Inf.\ Theory}, 22(6):644--654, 1976.

\bibitem{goldwasser1989knowledge}
S.~Goldwasser, S.~Micali, and C.~Rackoff,
``The Knowledge Complexity of Interactive Proof Systems,''
\emph{SIAM J.\ Comput.}, 18(1):186--208, 1989.

\bibitem{boneh2001bls}
D.~Boneh, B.~Lynn, and H.~Shacham,
``Short signatures from the Weil pairing,''
\emph{ASIACRYPT 2001}, pp.~514--532, 2001.

\bibitem{unger2015sok}
N.~Unger, S.~Dechand, J.~Bonneau, S.~Fahl, H.~Perl, I.~Goldberg, and
M.~Smith,
``SoK: Secure Messaging,''
\emph{IEEE S\&P}, pp.~232--249, 2015.

\bibitem{chase2019signal}
M.~Chase, T.~Perrin, and G.~Zaverucha,
``The Signal Private Group System,''
\emph{CCS 2020}, 2020.

\bibitem{blake2020post}
M.~Blake-Wilson, S.~Paquin, and D.~Stebila,
``Post-Quantum TLS,''
\emph{IETF Draft}, 2020.

\bibitem{langley2016cecpq2}
A.~Langley,
``CECPQ2: Post-quantum key agreement,''
\emph{Google Security Blog}, 2019.

\bibitem{dworkin2007nist}
M.~Dworkin,
``Recommendation for Block Cipher Modes: Galois/Counter Mode,''
\emph{NIST SP 800-38D}, 2007.

\bibitem{rescorla2018tls13}
E.~Rescorla,
``The Transport Layer Security (TLS) Protocol Version 1.3,''
\emph{RFC 8446}, 2018.

\bibitem{perrin2016double}
T.~Perrin and M.~Marlinspike,
``The Double Ratchet Algorithm,''
\emph{Signal Foundation}, 2016.

\bibitem{cremers2017automated}
C.~Cremers and M.~Feltz,
``Beyond eCK: Perfect Forward Secrecy under Actor Compromise and Ephemeral-Key Reveal,''
\emph{Des.\ Codes Cryptogr.}, 74:183--218, 2015.

\bibitem{nist2024pqc}
NIST,
``Post-Quantum Cryptography Standardization,''
\emph{csrc.nist.gov}, 2024.

\end{thebibliography}

\appendix

\section{PSP Message Wire Format}
\label{app:wire}

\begin{lstlisting}[caption={Binary Wire Format}]
PSP Pairwise Message (variable length):
  +--------+--------+-------+-------+
  | Version| Type   | DHPub | Prev  |
  | 1 byte | 1 byte |32 byte| 4 byte|
  +--------+--------+-------+-------+
  | Counter| Nonce  | Ciphertext    |
  | 4 byte |12 byte | variable      |
  +--------+--------+---------------+
  | AEAD Tag                        |
  | 16 bytes                        |
  +---------- ---------------------+

PSP Group Message (variable length):
  +--------+--------+--------+------+
  | Version| Type   | GroupID| Sender|
  | 1 byte | 1 byte |32 byte|32 byte|
  +--------+--------+--------+------+
  | Epoch  | Counter| Nonce  | CT   |
  | 8 byte | 8 byte |12 byte| var  |
  +--------+--------+--------+------+
  | AEAD Tag| Ed25519 Signature     |
  | 16 byte | 64 bytes              |
  +---------+-----------------------+
\end{lstlisting}

\section{Key Derivation Functions}
\label{app:kdf}

\begin{lstlisting}[language=Go,caption={KDF Implementations}]
func KDF_RK(
    rk [32]byte, dhOut [32]byte,
) (newRK, ck [32]byte) {
    prk := hkdf.Extract(sha256.New,
        dhOut[:], rk[:])
    r := hkdf.Expand(sha256.New,
        prk, []byte("PSP-RK"), 64)
    copy(newRK[:], r[:32])
    copy(ck[:], r[32:64])
    return
}

func KDF_CK(
    ck [32]byte,
) (newCK, mk [32]byte) {
    mk = hmacSHA256(ck[:], []byte{0x01})
    newCK = hmacSHA256(ck[:], []byte{0x02})
    return
}

func hmacSHA256(
    key, data []byte,
) [32]byte {
    h := hmac.New(sha256.New, key)
    h.Write(data)
    var out [32]byte
    copy(out[:], h.Sum(nil))
    return out
}
\end{lstlisting}

\end{document}
